\section{The integration pipeline}\label{note:pipeline}

\subsection{Adapters}

Each of the 54 source databases is reached through one adapter module
implementing a common protocol. An adapter yields a stream of records, each
carrying the source's native identifier, a list of property values, an optional
structure, a list of cross-database identifiers the source is authoritative for,
and a flag stating whether the record is experimental. A property value is a
tuple of canonical property name, value, unit and method string.

Two adapter flags modify the pipeline. \texttt{properties\_only} forbids the
source from contributing structures to the pool, which is how licence
restrictions are enforced in code rather than by convention.
\texttt{always\_hash} forces fingerprint matching for sources whose published
cross-references are not authoritative, that is, whose cited identifier names a
related entry rather than the same entry.

Each source passes through three resumable stages, leaving its artefacts on disk
so that a stage can be re-run without repeating the ones before it: adapt, which
writes the staged property rows and the structures needing fingerprints; hash,
which computes fingerprints in parallel; and match, which resolves records
against the match map and separates hits from records to be admitted.

\subsection{Structural fingerprints}

Structures are fingerprinted with the BAWL hasher~\cite{ramlaoui2025lemat} under
the protocol used for the original deduplicated merge, so that fingerprints
computed here are comparable with those already in the pool: parse the structure;
reject cells with more than 500 sites; reject disordered occupancies; reduce to
the primitive cell with pymatgen~\cite{ong2013python}; reject primitive cells with
more than 300 sites; then hash. The resulting string has the form
\texttt{<md5>\_<space group>\_<formula>}, and the space-group token is reused to
complete space-group coverage for the whole pool without a second symmetry
analysis. Failures are recorded per record with a reason and never abort a run.

\subsection{Two-path matching and admission}

Matching proceeds in two ordered paths against a match map holding
fingerprint-to-material and (source, native identifier)-to-material entries. A
record is first looked up by every authoritative identifier it carries. Records
still unmatched are looked up by fingerprint. A record matching nothing is
admitted as a new pool material: its structure is appended to the structure blob,
its metadata rows are written, and its fingerprint and identifiers are inserted
into the match map, which is what makes the enlarged substrate available to
subsequent sources. The chemical formula is never used as a match criterion.

Every match writes a row recording the material, the source, the native
identifier and the route taken, with \texttt{match\_via} equal to
\texttt{id:<database>}, \texttt{hash} or \texttt{origin}. The \texttt{origin}
route records that a source contains the material as its own entry. For a source
that admitted the structure this means it contributed it; for a properties-only
source it means only that the source holds a corresponding record. The
distinction is load-bearing for the licence position in
Note~\ref{note:licence}.

\subsection{Database schema}

Five tables extend the structure pool. \texttt{source\_defs} holds one row per
database with its licence and canonical URL. \texttt{property\_defs} holds the
canonical property registry. \texttt{material\_properties} holds one row per
value with its material, property, source, native identifier, numeric or text
value and method. \texttt{material\_source\_refs} holds the provenance references
described above. \texttt{material\_curves} holds one row per stored spectrum with
its type, format, source, method and an offset into a separate binary blob.

No value is ever overwritten or averaged. Several sources reporting the same
property for the same material coexist as separate rows, and the interface
exposes them as such. The scalar filter columns on the structure table, which
exist so that common queries need not join the property table, are backfilled
from \texttt{material\_properties} following a fixed source priority; existing
values are never overwritten, and the backfilled column is a convenience view
over a record that remains complete.

\section{The domain where structural fingerprints fail}\label{note:fingerprint}

A bonding-graph fingerprint identifies a crystal by its connectivity and its
primitive composition. For crystals that differ in either, this is exactly the
right invariant. Enumerated framework libraries violate the premise.

Hypothetical zeolites and hypothetical metal--organic frameworks are generated by
decorating a small number of parent topologies. Sibling entries therefore share a
bonding graph and a primitive composition, which is the entire content of the
fingerprint. Measured on the Deem zeolite set, the fingerprint path collapses
126{,}578 distinct frameworks onto 281 pool materials, one of which absorbs
27{,}007 of them. On the hypothetical metal--organic framework set, 46{,}891 of
135{,}818 hashable records collided with an earlier record and were dropped by
the standard pipeline, discarding 281{,}346 property values.

These are not duplicates, and treating them as such is worse than losing them.
Two hypothetical frameworks differing only in a functional group have their own
pore geometry and their own grand-canonical carbon dioxide isotherm; routing both
through one pool material would place several mutually contradictory porosities
on one structure, with nothing downstream able to detect it.

We treat the failure as a property of the source rather than of the individual
record. Sources whose fingerprints provably cannot separate their own entries are
declared fingerprint-unreliable, their fingerprint matches are rejected at load
time, and their records are admitted under a composite key pairing the
fingerprint with the cell volume per site rounded to two decimals. Volume per
site is invariant to the choice of unit cell and separates frameworks that a
graph hash cannot. Records that still collide on the composite key are dropped
rather than merged. Critically, the colliding fingerprints are withheld from the
match map, so that no later source can resolve onto an ambiguous structure.
Composite admission recovered 38{,}571 framework structures and 358{,}924
property values that the fingerprint path had discarded, and purged 28{,}914
ambiguous fingerprints that an earlier load had already registered.

The remedy is specific, not general. Any future source built by decorating a
fixed topology will need the same detection and the same treatment.

\section{Provenance and the canonical property registry}\label{note:provenance}

Property names are mapped onto a canonical registry, each entry with a fixed unit
and value kind. Mapping is not automatic: a source field enters the registry
under an existing name only when the quantity is the same physical quantity, and
otherwise gains its own name.

Two decisions show where the line falls. The residual of a geometric restraint
optimisation reported by a structure-prediction code is registered as
\texttt{structure\_prediction\_r\_factor} rather than \texttt{r\_factor}, because
the crystallographic refinement residual it would otherwise join has a different
meaning and an overlapping numeric range, so no downstream filter could separate
the two. Framework energies quoted per tetrahedral site in kJ\,mol$^{-1}$ are
registered as \texttt{energy\_per\_t\_atom} rather than folded into
\texttt{total\_energy\_per\_atom} in eV\,atom$^{-1}$.

Method strings are required to be distinguishing. Where one source reports the
same property for the same material under several settings, each value carries a
method string that separates them, and the loader treats (native identifier,
property, method) as the uniqueness key. How far this has to go is set by the
source rather than by the property. Band gap carries 12{,}275 distinct method
strings across the database, but 12{,}226 of them belong to the JARVIS tight-binding
dataset and take the form \texttt{QETB (tight-binding) @ 13.22 \AA$^{3}$/atom}:
its fitting set contains many deliberately non-equilibrium geometries of the same
compound, so the cell volume per atom is what separates one value from another
and the method string has to carry it. The remaining 49 strings are the
functional and setting labels of every other source combined. Without the volume
in the string those 12{,}226 values would collapse onto one per material, and the
fidelity gap between tight-binding and density functional theory visible in
Fig.~2f of the main text would be unrecoverable.

\section{Completeness audit}\label{note:audit}

The audit is deliberately independent of the loader. For each source it requires
three separate equalities to hold, each able to fail on its own.

\textbf{Load integrity.} The row count the loader reported must equal the row
count the database holds for that source. A mismatch means a load was interrupted,
ran twice, or was cleared and never re-inserted. A summary table alone would
never show this, because it reports only what the database currently contains.

\textbf{Staging yield.} Every property row the adapter staged must be either
present in the database or accounted for as a native identifier that no pool
material corresponds to. Staged minus loaded minus unmatched must be zero;
anything else is silently dropped data.

\textbf{Field coverage.} Every distinct property name present in the staged data
must exist in the database for that source. This is the check that catches the
failure mode that matters most: a field the adapter maps but the loader never
lands, which leaves row totals plausible while an entire quantity goes missing.

Sources staged on disk but absent from the source table are reported as never
loaded, and sources in the registry with no staging directory are reported as not
staged, so nothing can escape scrutiny by not being mentioned. Applied to the
completed database, the audit reports 54 sources and no discrepancies.

Each check is computed with a single grouped scan over the property table. This
is not a stylistic preference. The natural formulation, a correlated subquery
counting rows per source, rescans 121 million rows once per source, takes about
an hour, and evaluates each source against a different snapshot of a database
other processes may be writing. Every aggregate statistic reported in this work
uses grouped scans for the same reason.

\section{Licence handling and the ICSD position}\label{note:licence}

Licences are enforced at three points. Sources that may not redistribute
structures are marked properties-only in code, so no structure of theirs can
enter the pool even if a future adapter change would otherwise admit one. Sources
whose terms prohibit redistribution are excluded from any public snapshot. And
the distinction between an \texttt{origin} reference as contribution and an
\texttt{origin} reference as correspondence is maintained, because it is the
difference between holding a licensed structure and holding a pointer to one.

The Inorganic Crystal Structure Database is the case where this matters most. Its
adapter is properties-only. It annotates pool materials with chemical and mineral
names and calculated densities, and its 73{,}823 origin references are provenance
pointers marking pool materials that correspond to an ICSD collection code. Its
count of admitted structures is zero, its staged structure file is empty, and no
pool material lists it among its contributing sources. No FIZ Karlsruhe structure is
stored or redistributed by this work.

Four sources restrict redistribution in total, contributing 31{,}642{,}922
values, or 26.0\% of the database. They are searchable in a database built locally
from the published pipeline and are filtered out before public release. Any
statistic quoted from this work therefore has to state whether it describes the
locally built database or the public snapshot; every number in the main text
describes the former.

\section{Ingestion order and re-entrancy}\label{note:order}

Because each admitted structure enlarges the match map, a source ingested early
is matched against a smaller substrate than the same source ingested late. Ingestion
order is therefore not neutral, and Table~\ref{tab:sources} should be read with
the load order of Fig.~2c in mind.

The pipeline is built to be re-entrant so that this is recoverable: re-running a
source after the pool has grown is additive rather than destructive, because the
loader keys property rows on (native identifier, property, method) and reference
rows on (material, source, native identifier), and admitted fingerprints are
already registered. Two consequences follow for anyone reproducing this work.
Re-running an early source will attach references that the first pass could not
resolve, and the absolute figures will therefore rise rather than change
character. And a database that appears to contribute nothing new may simply have
been ingested before the structures it describes had arrived.

\section{Cross-database agreement protocol}\label{note:agreement}

Agreement was computed for eight properties whose values are directly comparable
between databases once the method is held fixed. Seven have at least one
structure described by two databases and are reported; volume per atom is
supplied by a single source and drops out. Extensive and convention-dependent
quantities were excluded by construction, in particular total energies and
prototype labels: two databases disagreeing on a total energy reflects a
different reference state rather than a different answer.

For each property, all values were read in one sequential pass and grouped by
material. Within a material each source was reduced to the median of its own
values, so that a source reporting a quantity under several methods is not
counted more than once, and the spread was taken as the range across sources,
absolute for energies, magnetisation and moduli and relative for density. The
tolerance within which two values are called concordant is 50~meV for band gap,
50~meV\,atom$^{-1}$ for formation energy and hull distance, 0.05~$\mu_{\rm B}$
for magnetisation, 5~GPa for the moduli and 2\% for density. Source pairs sharing
fewer than 200 materials are not reported. The complete pair table is
Table~\ref{tab:pairs}.

Two cautions apply to its interpretation. Agreement between two databases is not
evidence of independence: pairs whose median absolute difference is exactly zero
are carrying the same underlying calculation, and the table separates them from
genuinely independent pairs by that signature rather than by provenance metadata,
which the sources do not always supply. And a tolerance is a convention; the
distributions in Fig.~2e are given in full so that a reader preferring a different
threshold can read off the corresponding figure.

\section{Cost-tier assignment}\label{note:tiers}

The 69 general-purpose properties of Fig.~3b were assigned by hand to four tiers
according to the calculation they require: readable from a single self-consistent
run; requiring a small number of extra self-consistent runs such as a slab or a
strained cell; requiring linear response or finite differences over many
displacements; or requiring anharmonic, transport or beyond-semilocal treatment.

Properties whose domain is one material family, such as gas uptake or mineral
name, were excluded, because their coverage measures the size of that family
rather than the price of the calculation. The assignment is a judgement and is
listed in full in Table~\ref{tab:tiers} so that it can be disputed item by item.
The conclusion drawn from it in the main text, that no property requiring a
calculation of its own reaches a tenth of the pool, does not depend on the tier
boundaries: it depends only on the separation between properties that fall out of
a calculation already being performed and properties that do not.

\section{Census and construction of the training sets}\label{note:census}

\subsection{What counts as one task}

A task is one property at one level of theory, not one property. Method strings
are collapsed onto a level of theory by a fixed mapping, so that semilocal,
PBEsol, SCAN, r$^2$SCAN, hybrid, meta-GGA, tight-binding and experimental values
filed under a single property name become separate targets. Under that definition
the 151 populated properties yield 497 candidate tasks. Counting the tasks with
at least $10^6$, $10^5$, $10^4$, $2\times10^3$ and $3\times10^2$ trainable
crystals gives 13, 97, 161, 210 and 249 regression tasks respectively, over 13,
47, 77, 97 and 110 distinct properties, and 1, 15, 31, 37 and 40 classification
tasks. Every structure in the pool has a stored geometry, so no task has labels
without inputs; candidate structures average 15.1 sites and are capped at 500.

\subsection{Four constraints the corpus imposes on the labels}

\textbf{Mixed levels of theory.} Of the 7.82 million crystals carrying a band
gap, 81\% are semilocal, 8\% tight-binding and 4\% each PBEsol and SCAN, with
HSE06, meta-GGA, GLLB-SC and measured values making up 13 levels in total.
Offsets between levels are of order 1~eV and cannot be inferred from structure,
so a model trained on the pooled column learns which database supplied the row.
The default policy keeps the largest level per property; a split policy trains
one model per level.

\textbf{Condition-dependent quantities.} Some properties carry several values per
crystal that are all correct. Carbon dioxide uptake is reported at four
grand-canonical conditions, hydrogen uptake at five temperature and pressure
combinations, decomposition energy against four different reference hulls,
migration barrier by method, dimensionality and working ion, and topological
class with and without spin--orbit coupling. Enumerated zeolite libraries store
directional components of the Poisson ratio and stress-axis components of the
Young's modulus rather than an isotropic average. Taking a median over these
yields a number that is not any physical quantity, so each condition becomes its
own task.

\textbf{Cross-database disagreement.} Across the 162 constructed sets,
5{,}831{,}192 crystals carry labels from two or more databases and 519{,}599 of
them disagree by more than a quarter of the target's standard deviation. The
resolved label is the median, but the spread and the number of contributing
sources are written alongside it, so that the disagreement can weight the loss
rather than disappear into it.

\textbf{Split leakage.} The same compound recurs under different accession
numbers, so a random split places identical chemistry on both sides and inflates
the reported score in exactly the regime a completion is wanted for. Holding out
whole chemical systems is the default, but it degenerates where a task occupies a
single system: 135{,}887 hypothetical zeolites are all SiO$_2$ and the carbon
allotrope library is all carbon, so a chemistry-level holdout leaves both the
validation and the test fold empty. This happened for 23 of the 162 sets before
the rule was made to fall back through formula to random and to record the level
actually used. Splits are cut by material count rather than by group, because
hashing groups directly sends 29\% of the zeolites into a fold targeted at 10\%.

\subsection{Applicability domains}

Each task carries two descriptions of the structures it may be asked about. The
\emph{declared} domain is written by hand and states what the property requires
of a structure: 78 tasks are pool-wide, 25 require accessible porosity, 22 a
mobile working ion, 15 a gap, 4 a framework topology, 2 a magnetically ordered
ground state, 2 a van der Waals gap and 2 are surface quantities held against
bulk cells. Five are labelled synthesized-only, where a missing label means
nobody measured it rather than that the quantity is undefined. Seven are labelled
biased-labels: geometric pore descriptors are computable for any periodic
structure and give a well-defined answer, usually near zero, for a dense one, so
their concentration in porous materials is a sampling bias rather than a domain
restriction, and the two are handled differently.

The \emph{empirical} domain is measured from the training set: the family
distribution, the number of chemical systems, the 99th percentiles of density and
site count, and the elements present in every training structure. The last is the
hardest constraint. Carbon dioxide uptake has never been trained on a
carbon-free material; lithium migration barriers have never been trained on a
lithium-free one. Disagreement between the declared and the empirical domain is
itself informative: a property declared pool-wide whose labels sit in one family
is biased rather than restricted.

Enforcing the domains, 35 tasks may be asked about at least 90\% of the pool, 22
about 50 to 90\%, 36 about 10 to 50\%, 39 about 1 to 10\% and 30 about less than
1\%. The total number of admissible predictions is 871{,}097{,}089 against
2{,}194{,}010{,}874 for 162 tasks over the whole pool. One case is worth
recording: a force-field bulk modulus trained on 135{,}817 of the pool's
136{,}305 zeolites has 474 structures left to predict inside its validated
domain, and its 13.4 million apparent gaps are all outside it.

The gate is structural rather than geometric. Porosity is enforced through family
labels rather than through a computed pore-limiting diameter compared against the
kinetic diameter of the adsorbate, which is the correct test and requires a
geometry analysis over the whole pool that has not been run. A permissive mode
that admits unlabelled structures whose density falls in the training band is
available for candidate generation only, and has a known failure direction, since
lithium metal passes a porosity density gate without having pores. Seventeen
tasks are gated on a gap or a magnetic moment that the database does not hold for
75{,}674{,}942 structure-task pairs; those decisions are deferred until the
corresponding tier-1 models have run.
